\documentclass[11pt]{article}
\usepackage[margin=1in]{geometry}
\usepackage{xcolor}
\usepackage{tcolorbox}
\usepackage{hyperref}
\usepackage{enumitem}
\usepackage{booktabs}

% Define colors
\definecolor{osaorange}{RGB}{220,115,38}
\definecolor{aiblue}{RGB}{30,144,255}

% Configure hyperref
\hypersetup{
    colorlinks=true,
    linkcolor=blue,
    urlcolor=blue,
    citecolor=blue
}

\title{\textbf{H524 Introduction to Biostatistics\\
\Large Mid-Semester Individual Assignment\\
AI-Assisted Statistical Analysis\\
Fall 2025}}
\author{Dr. John Molitor}
\date{Due: Week 6 (Friday by 11:59 PM)}

\begin{document}
\maketitle

\begin{tcolorbox}[colback=aiblue!10,colframe=aiblue,boxrule=2pt,arc=2pt,left=6pt,right=6pt,top=6pt,bottom=6pt]
\centering
\textbf{\large Assignment Overview}\\[0.3em]
This individual assignment replaces the traditional midterm exam and serves as preparation for your group project. You will conduct a focused statistical analysis using AI tools, demonstrating your understanding of descriptive statistics, hypothesis testing, and AI integration principles covered in the first half of the course.
\end{tcolorbox}

\section{Learning Objectives}

By completing this assignment, you will demonstrate your ability to:
\begin{itemize}[leftmargin=20pt]
\item Apply descriptive statistical methods to real health data
\item Conduct appropriate hypothesis tests and interpret results
\item Effectively use AI tools for statistical analysis while maintaining scientific rigor
\item Document and verify AI-generated outputs
\item Communicate statistical findings clearly and accurately
\end{itemize}

\section{Assignment Components}

\subsection{Dataset Selection}
Choose \textbf{ONE} of the following simplified datasets:

\subsubsection{Option 1: Heart Disease Risk Factors}
\begin{itemize}
\item \textbf{Data:} Framingham Heart Study subset (provided on Canvas)
\item \textbf{Variables:} Age, gender, BMI, systolic BP, cholesterol, smoking status, diabetes status
\item \textbf{Sample Size:} ~500 participants
\item \textbf{Focus:} Cardiovascular risk factor analysis
\end{itemize}

\subsubsection{Option 2: COVID-19 Vaccination Outcomes}
\begin{itemize}
\item \textbf{Data:} Simulated vaccination study data (provided on Canvas)
\item \textbf{Variables:} Age group, vaccination status, breakthrough infections, hospitalization
\item \textbf{Sample Size:} ~1000 participants
\item \textbf{Focus:} Vaccine effectiveness analysis
\end{itemize}

\subsubsection{Option 3: Mental Health and Exercise}
\begin{itemize}
\item \textbf{Data:} Depression screening study (provided on Canvas)
\item \textbf{Variables:} Age, gender, exercise frequency, depression score, anxiety score
\item \textbf{Sample Size:} ~400 participants
\item \textbf{Focus:} Mental health and lifestyle factors
\end{itemize}

\section{Required Analysis Tasks}

\subsection{Task 1: Descriptive Analysis (25 points)}
\begin{enumerate}
\item \textbf{Data Exploration:}
   \begin{itemize}
   \item Import and examine the dataset structure
   \item Check for missing values and outliers
   \item Create summary statistics for all variables
   \end{itemize}

\item \textbf{Descriptive Statistics:}
   \begin{itemize}
   \item Calculate measures of central tendency and variability for continuous variables
   \item Create frequency tables for categorical variables
   \item Generate appropriate summary tables
   \end{itemize}

\item \textbf{Data Visualization:}
   \begin{itemize}
   \item Create 3-4 appropriate graphs (histograms, box plots, bar charts)
   \item Include proper titles, labels, and legends
   \item Briefly interpret each visualization
   \end{itemize}
\end{enumerate}

\subsection{Task 2: Hypothesis Testing (30 points)}
\begin{enumerate}
\item \textbf{Research Question:}
   \begin{itemize}
   \item Formulate one clear, testable research question based on your chosen dataset
   \item State your null and alternative hypotheses
   \end{itemize}

\item \textbf{Statistical Test Selection:}
   \begin{itemize}
   \item Choose the appropriate statistical test (t-test, chi-square, ANOVA)
   \item Justify your choice based on variable types and research question
   \item Check test assumptions
   \end{itemize}

\item \textbf{Analysis and Interpretation:}
   \begin{itemize}
   \item Conduct the statistical test
   \item Report test statistic, p-value, and confidence interval (if appropriate)
   \item Interpret results in the context of your research question
   \item Discuss practical significance vs. statistical significance
   \end{itemize}
\end{enumerate}

\subsection{Task 3: Written Report (20 points)}
\begin{enumerate}
\item \textbf{Report Structure (3-4 pages):}
   \begin{itemize}
   \item Introduction with research question and dataset description
   \item Methods section describing your analytical approach
   \item Results section with tables, figures, and statistical findings
   \item Brief discussion of findings and limitations
   \item References (course materials and dataset sources)
   \end{itemize}

\item \textbf{Communication Quality:}
   \begin{itemize}
   \item Clear, professional writing appropriate for a public health audience
   \item Proper statistical terminology and interpretation
   \item Well-formatted tables and figures with informative captions
   \end{itemize}
\end{enumerate}

\section{Deliverables}

Submit the following files via Canvas by the deadline:

\begin{enumerate}
\item \textbf{R Script File (.R):}
   \begin{itemize}
   \item Well-commented R code for all analyses
   \item Clear section headers for each task
   \item Working code that reproduces all results
   \end{itemize}

\item \textbf{Written Report (.pdf):}
   \begin{itemize}
   \item 3-4 page report following the specified structure
   \item Include all tables and figures within the report
   \item Professional formatting and presentation
   \end{itemize}

\item \textbf{Data File (.csv):}
   \begin{itemize}
   \item Your cleaned dataset (if modifications were made)
   \item Include variable labels and coding information
   \end{itemize}
\end{enumerate}

\section{Grading Rubric}

\begin{center}
\begin{tabular}{lcc}
\toprule
\textbf{Component} & \textbf{Points} & \textbf{Criteria} \\
\midrule
Descriptive Analysis & 35 & Comprehensive data exploration and visualization \\
Hypothesis Testing & 40 & Appropriate test selection and correct interpretation \\
Written Report & 25 & Clear communication and professional presentation \\
\midrule
\textbf{Total} & \textbf{100} & \\
\bottomrule
\end{tabular}
\end{center}

\subsection{Detailed Grading Criteria}

\subsubsection{Excellent (90-100\%)}
\begin{itemize}
\item Comprehensive and accurate statistical analysis
\item Clear, professional writing with correct interpretations
\item Well-designed visualizations and properly formatted tables
\end{itemize}

\subsubsection{Good (80-89\%)}
\begin{itemize}
\item Mostly correct analysis with minor computational errors
\item Generally clear writing with mostly correct interpretations
\item Appropriate visualizations with minor formatting issues
\end{itemize}

\subsubsection{Satisfactory (70-79\%)}
\begin{itemize}
\item Basic analysis completed with some errors
\item Unclear writing or some misinterpretations
\item Basic visualizations that convey main findings
\end{itemize}

\subsubsection{Needs Improvement (<70\%)}
\begin{itemize}
\item Major errors in statistical analysis or interpretation
\item Poor communication or significant misunderstandings
\item Inadequate or missing visualizations
\end{itemize}

\section{Resources and Support}

\subsection{Statistical Resources}
\begin{itemize}
\item Course lecture notes and textbook (Pagano \& Gauvreau)
\item R documentation and help files
\item Canvas discussion board for questions
\end{itemize}

\subsection{AI Tool Access}
\begin{itemize}
\item GitHub Copilot (required - set up in RStudio)
\item ChatGPT, Gemini, or similar AI assistant for consultation
\item Course-provided AI prompt templates
\end{itemize}

\subsection{Technical Requirements}
\begin{itemize}
\item R and RStudio (latest versions)
\item Required R packages: tidyverse, ggplot2, stats
\item PDF generation capability for reports
\end{itemize}

\section{Academic Integrity}

\begin{tcolorbox}[colback=red!5,colframe=red,boxrule=1pt,arc=2pt,left=6pt,right=6pt,top=6pt,bottom=6pt]
\textbf{Important Reminders:}
\begin{itemize}
\item This is an \textbf{individual assignment} - no collaboration with classmates
\item All AI-generated content must be verified and properly documented
\item You are responsible for the accuracy of all submitted work
\item Failure to document AI usage constitutes academic dishonesty
\item Submitting unverified AI output as your own work is plagiarism
\end{itemize}
\end{tcolorbox}

\section{Tips for Success}

\begin{enumerate}
\item \textbf{Start Early:} Data analysis often takes longer than expected
\item \textbf{Document Everything:} Keep detailed notes of your analysis decisions
\item \textbf{Verify AI Output:} Always double-check AI-generated code and interpretations
\item \textbf{Ask Questions:} Reach out if you encounter difficulties
\item \textbf{Focus on Understanding:} Ensure you can explain every step of your analysis
\item \textbf{Practice Presentation:} You may be asked to briefly present your findings in class
\end{enumerate}

\section{Late Submission Policy}

\begin{itemize}
\item Late submissions will be penalized 10\% per day
\item Extensions may be granted for documented emergencies
\item Technical difficulties are not grounds for extensions - submit early to avoid issues
\end{itemize}

\end{document}